% ARCH_ALG task 3. New IDs ARCH-15-*. No Masullo source material used.
\section{Cote 15: actual forward cycles and the precise mixed Tate cases}
\label{arch15-core}

\paragraph{Archive locator and scope.}
The question concerns cote 15, physical PDF pp.~94--106: the question
\(Q'\) is on physical p.~94, footer 93,
\texttt{ARCHIVE\_15\_scan094.tex}; the norm discussion is on physical
p.~98, footer 97, \texttt{ARCHIVE\_15\_scan098.tex}; and the boxed purity
statement and conjecture forte are on physical p.~100, footer 99,
\texttt{ARCHIVE\_15\_scan100.tex}.  The spectral and finite-certificate
results of checkpoint 09 are inputs, not new theorems here. In particular,
CC-XC03 is the general projected Tate criterion, CC-XC11--12 treat the
rational spectral target, CC-XC15 is the forward certificate, and CC-XC08
already includes curves and their products. The statements below identify
the remaining cycle spaces for the first two requested classes and give
actual geometric constructions for specified subclasses.

\paragraph{Coefficients and the input ledger.}
All Chow groups and Chow motives in this section have rational
coefficients. Put \(k=\mathbf F_q\), and use geometric Frobenius on
cohomology. Write \(\mathbb L=\mathbf Q(-1)\), of Frobenius eigenvalue
\(q\). A degree-zero correspondence from \(X\) to \(Y\) lies in
\(\mathrm{CH}^{\dim X}(X\times Y)_{\mathbf Q}\).
We distinguish two uses of an abelian target.

\begin{enumerate}
\item A \emph{geometrically supplied target} is a direct summand of a
finite sum of twisted abelian and Artin Chow motives, with its projector
and the forward map supplied by algebraic cycles.
\item The \emph{fixed spectral target} \(M(X)\) is the rational numerical
object of CC-XC11, specified by the Frobenius polynomials. Its general
existence and full realization in the checkpoint use the all-abelian Tate
input. Denote this hypothesis by \(\mathrm{AB}\).
\end{enumerate}

In a presentation of the second kind, \(e_N\) below is a numerical
idempotent, represented by rational Chow correspondences. Under
\(\mathrm{AB}\) it is idempotent in the homological category. We do
\emph{not} assert that the chosen Chow representative is idempotent
modulo rational equivalence; no nilpotence or projector-lifting theorem
is being silently used. Composition with that representative still
gives an actual Chow cycle and the required homological map. By
contrast the projectors \(ab\) in the explicit geometric constructions
are proved idempotent already in the Chow category.

In the checkpoint's dependency chain, \(\mathrm{AB}\) uses the named
outside preprint of van Hoften--Sempliner, \emph{The Hodge Conjecture for
CM abelian varieties}, Hexagon 2610.00259v2, followed by Milne's published
1999 theorem. It is not an unconditional published theorem in this
return. The other two named outside inputs, \emph{Milne's rationality
conjecture}, Hexagon 2610.00226v1, and Milne's note 2611.00244v1, are not
used in the arguments below. Neither A nor M is used anywhere in this
section. Formal deductions and displayed intersection calculations have
tag ``none''; the geometric and cohomological inputs are identified below
as published theorems. Inherited checkpoint assertions retain their
original dependencies when invoked.

\subsection{ARCH-15-CUBIC-GATE: the exact case for an arbitrary cubic fourfold}

Let \(X\subset\mathbf P^5_k\) be a smooth cubic fourfold, and let
\(H\in\mathrm{CH}^1(X)\) be its hyperplane class. The five rational Chow
projectors
\[
 p_{2i}=\frac13 H^{4-i}\times H^i\quad(0\leq i\leq4),\qquad
 p_{\mathrm{prim}}=\Delta_X-\sum_{i=0}^4p_{2i}
 \tag{ARCH-15.1}
\]
are mutually orthogonal, and \(p_{\mathrm{prim}}\) is self-transpose.
Its realization is exactly \(H^4_{\mathrm{prim}}(X)\), of dimension 22;
the complement is \(\bigoplus_{i=0}^4\mathbb L^i\).

Assume \(\mathrm{AB}\), and choose a splitting
\[
 M(X)\simeq\left(\bigoplus_{i=0}^4\mathbb L^i\right)\oplus N_X,
 \qquad N_X=(T,e_N),\quad T=\bigoplus_{\nu=1}^s h(Y_\nu),
 \tag{ARCH-15.2}
\]
where the \(Y_\nu\) are abelian varieties or finite etale schemes, as in
CC-XC11. The splitting chooses one of the hyperplane Tate lines inside
the middle Tate isotypic part; it does not require that part to have
multiplicity one. For \(\ell\ne p\), define
\[
 \begin{split}
 \kappa_{X,\ell}:\quad
 \bigoplus_\nu\mathrm{CH}^4(X\times Y_\nu)_{\mathbf Q}
       \otimes_{\mathbf Q}\mathbf Q_\ell
 &\longrightarrow
 \operatorname{Hom}_{\mathbf Q_\ell[F]}
       (H^4_{\mathrm{prim}}(X),R_\ell N_X),\\
 z&\longmapsto R_\ell(e_N z p_{\mathrm{prim}}).
 \end{split}
 \tag{ARCH-15.3}
\]

\textbf{Statement.} Under \(\mathrm{AB}\), the following are equivalent:
(i) actual rational Chow cross-cycles give a forward isomorphism from
\(h(X)\) to the fixed target \(M(X)\) in homological realization;
(ii) \(\kappa_{X,\ell}\) is onto for one \(\ell\ne p\);
(iii) for the displayed finite family, the usual Tate cycle maps in
codimension 4 on all \(X\times Y_\nu\) are onto. Once these conditions
hold, the rational certificate can be chosen to work for every
\(\ell\ne p\), as in CC-XC12 and CC-XC15. Condition (ii) is the smallest
projected formulation asserted here; condition (iii) is the equivalent
unprojected finite family under \(\mathrm{AB}\).

\textbf{Proof.} For product correspondences in (ARCH-15.1), composition
integrates the middle factor. Its degree is zero unless the middle
codimension is 4, and \(\deg H^4=3\). Thus
\(p_{2j}p_{2i}=\delta_{ij}p_{2i}\) in the rational Chow ring.
Transposition exchanges \(p_{2i}\) and \(p_{8-2i}\). Weak Lefschetz,
Poincare duality, and the middle Betti number 23 give the stated
realizations. Frobenius on \(H^4(X)\) is semisimple: this is the
Deligne--Rapoport Kuga--Satake result, explicitly recalled without a
characteristic restriction in Zarhin's published paper on ordinary
cubic fourfolds: in the inspected author preprint, introduction p.~1,
the sentence following the Betti-number calculation concerns an
arbitrary smooth cubic fourfold over a finite field of characteristic
\(p\), before ordinarity is imposed. Equivalently, the primitive
representation embeds, after a finite extension, into a tensor
construction on \(H^1\) of an abelian variety. The modern
characteristic-2 extension of the orthogonal Kuga--Satake map is also
described in Auel et al., proof of Theorem 4.10, author-hosted journal
PDF p.~20. Semisimplicity of a power
of an invertible characteristic-zero operator implies semisimplicity of
the operator itself.

Apply the inherited CC-XC03 argument to the self-dual summand
\((X,p_{\mathrm{prim}})\) and \(N_X\). The complement in (ARCH-15.1)
already has explicit algebraic maps to its Tate target. If an
isomorphism of the whole realization is initially given, transport
\(p_{\mathrm{prim}}\) to the abelian target; fullness there identifies
its image with the chosen \(N_X\). Thus no unjustified restriction of
an arbitrary isomorphism to a preselected summand is needed. This proves
(i)\(\Leftrightarrow\)(ii). Surjectivity of the unprojected Tate maps
implies (ii) by Kunneth and Poincare duality: the relevant component of
\(H^8(X\times Y_\nu,\mathbf Q_\ell(4))^{F=1}\) is the space of
degree-preserving Frobenius maps out of \(H^4(X)\). Conversely (i),
\(\mathrm{AB}\), and CC-XC04 imply Tate surjectivity on every
\(X\times Y_\nu\), giving (iii). Rationality and the simultaneous
prime-independent certificate are precisely the inherited CC-XC12--15
step. This proves the statement. \(\square\)

\paragraph{What the published ordinary Tate theorems do.}
Tate in codimension 2 on \(X\) itself is published in odd characteristic
by Madapusi Pera (2015, \S5.14 in the inspected arXiv v3) and in
characteristic 2 by Auel--Kulkarni--Petok--Weinbaum (Theorem 4.10,
\S4.6). Zarhin's Corollary 5.2 proves Tate for every self-power of an
\emph{ordinary} cubic fourfold. None of those statements asserts the
mixed codimension-4 algebraicity in (ARCH-15.3). Applying one of them as
though it already furnished that cross-cycle would leave a gap.

\subsection{ARCH-15-SPLIT-CUBIC: an explicit Chow construction}

Suppose \(p\ne3\) and
\[
 X=\{f(x_0,x_1,x_2)+g(x_3,x_4,x_5)=0\}\subset\mathbf P^5,
 \tag{ARCH-15.4}
\]
where \(E_f=\{f=0\}\subset\mathbf P^2\) and
\(E_g=\{g=0\}\subset\mathbf P^2\) are smooth plane cubics. Put
\[
 S_f=\{f(x)+u^3=0\}\subset\mathbf P^3,
 \quad S_g=\{g(y)+v^3=0\}\subset\mathbf P^3,
 \quad V=\operatorname{Bl}_{E_f\times E_g}(S_f\times S_g).
\]
The actual morphism
\[
 \phi:V\longrightarrow X,
 \qquad ((x:u),(y:v))\longmapsto(xv:-uy)
 \tag{ARCH-15.5}
\]
has degree 3. Consequently \(h(X)\) is a rational Chow direct summand
of \(h(V)\), with forward map \(a=\phi^*\) and reverse map
\(b=\frac13\phi_*\). After a finite extension \(k'/k\), there is an
explicit Chow isomorphism
\[
 h(V_{k'})\simeq
 (\mathbf1\oplus\mathbb L^{\oplus7}\oplus\mathbb L^2)^{\otimes2}
 \ \oplus\ h((E_f\times E_g)_{k'})\otimes\mathbb L.
 \tag{ARCH-15.6}
\]
Thus (ARCH-15.5), the blowup correspondences, and the rational-surface
correspondences give actual forward cycles into an abelian--Artin Chow
target over the original field.

\textbf{Proof and projector.} The smoothness, base ideal \((u,v)\),
resolution by the indicated blowup, and degree are the case
\(m=3,r=s=2\) of the already proved CC-XC10 join calculation. Here
\(\epsilon=-1\) already satisfies \(\epsilon^3=-1\); no roots of unity
extension is required for the map. The input geometry is the published
Shioda--Katsura Remark 1.10, p.~102. We use that inherited calculation,
not a second purported new proof of it.

For any proper generically finite morphism between smooth projective
varieties of equal dimension, graph composition gives
\[
 \phi_*\phi^*=3\Delta_X\quad\hbox{in }
 \mathrm{CH}^4(X\times X)_{\mathbf Q}.
 \tag{ARCH-15.7}
\]
Indeed the refined intersection giving this composition is the
fundamental class of the graph with its two equal images in \(X\);
pushing it to \(X\times X\) gives
\(\Delta_{X*}\phi_*[V]=3[\Delta_X]\). This also follows from the
projection formula after arbitrary smooth projective base change and
Manin's identity principle. In particular, this is a Chow identity,
not just an equality on cohomology. If \(\theta\) denotes (ARCH-15.6),
the target projector is
\[
 e=(\theta a)(b\theta^{-1}),\qquad
 e^2=e,\qquad (b\theta^{-1})(\theta a)=1_{h(X_{k'})}.
 \tag{ARCH-15.8}
\]

The smooth cubic surfaces become blowups of \(\mathbf P^2\) in six
points after a finite extension. The classical rational-surface
description and the Chow blowup formula give
\(h(S_f)=h(S_g)=\mathbf1\oplus\mathbb L^{\oplus7}\oplus\mathbb L^2\)
there. The center has codimension two, so its blowup contribution is
exactly one copy of \(h(E_f\times E_g)\otimes\mathbb L\); this proves
(ARCH-15.6). For clarity, with exceptional inclusion \(j\), bundle
projection \(\rho\), and blowdown \(\beta\), the blowup maps are
\((u,v)\mapsto\beta^*u+j_*\rho^*v\) and
\(w\mapsto(\beta_*w,-\rho_*j^*w)\). The minus sign follows from the
normal bundle \(\mathcal O(-1)\). The Chow blowup formula, or its
proof by the projective bundle and excess-intersection formulas, shows
these are inverse on Chow groups after every smooth projective product,
and hence inverse as Chow motives. A precise published locator for
the Chow formula is Jiang (2023), Theorem 2.8(3), p.~1472; the
motivic form is Corollary 3.10, p.~1484, with its stated blowup
specialization. Genus-one curves acquire origins
after a finite extension and then are elliptic curves.

The finite-field descent in CC-XC16 is also an actual rational Chow
construction. Induce the correspondences over \(k'\) by viewing their
schemes over \(k\); use restriction of scalars for their abelian
factors. For \(\tau:X_{k'}\to X\),
\(\tau_*\tau^*=[k':k]\Delta_X\); thus \(\tau^*\) and
\([k':k]^{-1}\tau_*\) give the final retraction. The induced abelian
factor is a summand of
\(h(\operatorname{Res}_{k'/k}A)\otimes h(\operatorname{Spec}k')\),
not necessarily of a single abelian variety without the Artin factor.
All denominators are permitted in rational Chow groups. The identities
descend because restriction followed by trace is multiplication by
\([k':k]\). This proves the assertion over \(k\). \(\square\)

\paragraph{Published-only scope.}
This is an application of CC-XC10, not a novelty claim. The actual Chow
construction uses published geometry and formal intersection theory,
and no A, M, or outside preprint. Moreover all abelian factors involved
are, geometrically, products of elliptic curves. Spie\ss's published
theorem and Milne (2001), Appendix A.7, first example on author p.~21,
prove Tate for \emph{every product and every power} of such elliptic
curves. This is the scope needed for endomorphisms, duality, projectors,
and mixed targets within this sector, rather than merely Tate for
divisors on an individual elliptic curve. Trace averages a cycle over
a field extension to give its invariant class over the original finite
field. Semisimplicity of abelian Frobenius and the invariant Poincare
pairing then give numerical--homological equality in this sector, by
the same elementary argument as CC-XC04. Thus full realization for this
chosen target sector uses published theorems alone. Identifying the
result with an \emph{arbitrary preselected} presentation of the global
CC-XC11 target still retains \(\mathrm{AB}\) unless that presentation
is itself known to lie in this published sector.

\subsection{ARCH-15-K3-GATE: finite-height K3 surfaces}

Let \(X/k\) be a K3 surface of finite height, in any characteristic.
Choose \(o\in\mathrm{CH}^2(X)_{\mathbf Q}\) of degree one; a closed
point divided by its degree suffices. Over a finite extension splitting
the geometric Neron--Severi group choose divisor classes
\(D_1,\ldots,D_\rho\) forming a rational basis, and put
\(B_{ij}=\deg(D_iD_j)\). Define
\[
 p_0=o\times X,\quad p_4=X\times o,\quad
 p_{\mathrm{NS}}=\sum_{i,j}(B^{-1})_{ij}D_i\times D_j,
 \quad p_{\mathrm{tr}}=\Delta_X-p_0-p_4-p_{\mathrm{NS}}.
 \tag{ARCH-15.9}
\]
The matrix is invertible by the Hodge index theorem. These are rational
Chow projectors; \(p_{\mathrm{NS}}\) and \(p_{\mathrm{tr}}\) descend to
\(k\). Let \(T_\ell(X)=\mathrm{NS}(X_{\bar k})_{\mathbf Q_\ell}^{\perp}
\subset H^2(X_{\bar k},\mathbf Q_\ell)\), of dimension \(22-\rho\).
Under \(\mathrm{AB}\), split from \(M(X)\) the known Artin--Tate
realization of \(p_0+p_4+p_{\mathrm{NS}}\); denote its complement by
\(N_X=(\bigoplus_\nu h(Y_\nu),e_N)\). The exact remaining map is
\[
 \bigoplus_\nu\mathrm{CH}^2(X\times Y_\nu)_{\mathbf Q}
       \otimes\mathbf Q_\ell
 \longrightarrow
 \operatorname{Hom}_{\mathbf Q_\ell[F]}(T_\ell(X),R_\ell N_X),
 \qquad z\longmapsto R_\ell(e_Nzp_{\mathrm{tr}}).
 \tag{ARCH-15.10}
\]

\textbf{Statement.} Under \(\mathrm{AB}\), forward rational Chow
cross-cycles to the fixed \(M(X)\) exist if and only if (ARCH-15.10) is
surjective for one \(\ell\ne p\). Equivalently, Tate in codimension 2
holds for the displayed finite family \(X\times Y_\nu\). The
conclusion is again simultaneous in \(\ell\) by CC-XC12--15.

\textbf{Proof.} The identities among (ARCH-15.9) follow by middle-factor
intersection, just as for (ARCH-15.1); \(p_{\mathrm{NS}}^2=p_{\mathrm{NS}}\)
is the matrix identity \(B^{-1}BB^{-1}=B^{-1}\). The other products
vanish by dimension and \(\deg o=1\). A K3 surface has
\(\operatorname{Pic}^0=0\), so equality of Neron--Severi divisor classes
is equality in \(\mathrm{CH}^1\). The tensor defined by the inverse
intersection form is invariant under change of basis and under Galois.
Consequently \(p_{\mathrm{NS}}\), initially over a finite Galois
extension \(k'/k\), has invariant restriction. Its trace divided by
\([k':k]\) restricts back to it. Restriction is injective on rational
Chow groups because trace composed with restriction is that same
nonzero integer; hence the projector identities and self-transposition
descend. This establishes the asserted Chow, rather than merely
cohomological, projector.

Frobenius on K3 cohomology is semisimple by the Kuga--Satake
cohomological embedding and semisimplicity for abelian varieties. The
all-characteristic construction and its integral-model corrections are
recorded by Ito--Ito--Koshikawa, \S1.3 and Remarks 6.9--6.10; their
\S10 uses this semisimplicity explicitly. The published Tate theorem
identifies geometric divisor classes with the root-of-unity part of
\(H^2(1)\), so \(T_\ell(X)(1)\) has no root-of-unity eigenvalues. The
known complement in (ARCH-15.9) is an Artin representation tensored
with \(\mathbb L\), together with \(\mathbf1\) and \(\mathbb L^2\),
and has actual algebraic maps to an Artin--Tate target by its divisor
basis and inverse intersection matrix. Apply CC-XC03 to the self-dual
remaining summand. The proof of the equivalence with unprojected Tate
and of rational simultaneous certification is the same application of
CC-XC04 and CC-XC12--15 as above. \(\square\)

\paragraph{The CM-lifting route and its precise stop.}
Ito--Ito--Koshikawa's published Theorem 1.1 gives a CM lifting after
finite extension for every finite-height K3 surface, and their Theorem
1.5 proves Tate for the square of every K3 surface, including
characteristics 2 and 3. The mixed Tate class in (ARCH-15.10) is not a
class on that square. A Kuga--Satake embedding into endomorphisms of
abelian cohomology is a cohomological construction; its algebraicity as
a cycle on a K3--abelian product is an additional assertion. Neither a
CM lifting nor the cited square theorem makes that assertion. The next
geometric route below supplies the missing cycle for a concrete
finite-height class.

\subsection{ARCH-15-KUMMER: actual degree-two Chow cycles}

Let \(p\ne2\), let \(A/k\) be an abelian surface, and put
\(K=\operatorname{Km}(A)\). Blowing up the finite etale scheme
\(A[2]\) gives \(\widetilde A\); the lifted involution has smooth Kummer
quotient, and the quotient morphism \(\phi:\widetilde A\to K\) has
degree two. Thus
\[
 h(\widetilde A)\simeq h(A)\oplus h(A[2])\otimes\mathbb L,
 \qquad
 a=\phi^*,\quad b=\tfrac12\phi_*,\quad ba=1_{h(K)}.
 \tag{ARCH-15.11}
\]
Composing with this blowup isomorphism and taking the projector \(ab\)
is an actual rational Chow construction of the forward target. It
requires no outside preprint, even for a general abelian surface.

\textbf{Proof.} The standard Kummer construction in characteristic
different from 2 is Huybrechts, \emph{Lectures on K3 Surfaces}, Chapter
1, Example 1.3(iii), author draft pp.~8--9. It supplies the smooth
quotient and degree. The center is smooth of codimension two and has
16 geometric points. Apply the same Chow blowup maps and graph
composition identity proved above, now with degree two. This yields
(ARCH-15.11) and the idempotent. If the construction or a chosen
abelian presentation is made over a finite extension, the rational
trace descent already checked above returns the cycles to \(k\).
\(\square\)

In particular, take \(A=E_1\times E_2\). If both elliptic curves are
ordinary, the slopes of \(H^2(K)\) are \(0,1^{20},2\); if exactly one
is ordinary they are \((1/2)^2,1^{18},(3/2)^2\). Indeed take the second
exterior power of the four \(H^1(A)\) slopes, and add the sixteen slope-1
exceptional classes in (ARCH-15.11). Pullback identifies the two
degree-2 cohomology spaces, since it is injective and both have
dimension 22. The K3 slope formula is the Newton
polygon with vertices \((0,0),(h,h-1),(22-h,21-h),(22,22)\),
together with the formal-Brauer interpretation of \(h\): Huybrechts,
Chapter 18, \S3.2--3.3, author draft pp.~399--401, especially the
polygon on p.~400 and (3.5) on p.~401. It gives height 1 and height 2 respectively. These are finite-height
examples in the requested second class. All powers and mixed products
of their chosen elliptic--Artin targets satisfy the published Tate
theorem cited above. The actual construction and full realization in
that sector therefore use published theorems alone. If both elliptic
curves are supersingular, the same construction is valid, but it falls
in the supersingular class already treated in CC-XC16--23.

\subsection{ARCH-15-CURVE-SQUARE: inherited construction and its scope}

Let \(C/k\) be smooth projective and geometrically connected of genus
\(g\ge2\), and let \(J=J(C)\). The answer to the construction question
for \(C\times C\) is already affirmative in CC-XC08(1), including its
finite-field descent. It is not a first new example of this return.
In the classical rational Chow notation the construction is
\[
 h(C)\simeq\mathbf1\oplus h^1(J)\oplus\mathbb L,
 \qquad
 h(C\times C)\simeq
 (\mathbf1\oplus h^1(J)\oplus\mathbb L)^{\otimes2}.
 \tag{ARCH-15.12}
\]
Here the curve projectors are \(o\times C\), \(C\times o\), and their
complement, for a degree-one rational zero-cycle \(o\); the middle
identification is the classical curve--Jacobian Chow correspondence.
Tensor the actual forward and reverse correspondences for the curve
and take their tensor product projectors. This proves the asserted
actual construction using the inherited CC-XC08 case, or the stronger
standard Chow form (ARCH-15.12), without A, M, or outside preprints.

Ordinary Tate on the surface \(C\times C\) is a published consequence
of Tate's theorem on homomorphisms of abelian varieties: its middle
mixed part after the Tate twist is \((H^1(C)\otimes H^1(C))(1)\), and divisor correspondences
identify its invariants with \(\operatorname{End}_{k}(J)\otimes
\mathbf Q_\ell\); the remaining invariant classes are products of
degree-one zero-cycle classes. This assertion about a surface is not
the all-codimension Tate assertion for arbitrary powers of its
Jacobian. Consequently, recognizing (ARCH-15.12) as the arbitrary
global spectral target retains \(\mathrm{AB}\) in general. If
\(J_{\bar k}\) is isogenous to a product of elliptic curves, the
published elliptic-product theorem does suffice for every power and
mixed product needed in this particular sector. This is an explicit
published-only sufficient subclass, not a claim that the general
Jacobian case is settled by the divisor theorem.

\subsection{ARCH-15-INPUT-AUDIT: what can and cannot replace the preprints}

\begin{center}
\begin{tabular}{p{0.35\linewidth}p{0.23\linewidth}p{0.33\linewidth}}
Step & Input tag & Exact scope\\\hline
Hyperplane and divisor projectors; trace; graph composition
 & none (formal calculation) & Actual rational Chow identities\\
Join geometry; rational surfaces; blowup; Kummer; curve--Jacobian
 & published theorem & Geometric Chow targets in the displayed classes\\
Tate on elliptic products, all powers included
 & published theorem & Full realization in the elliptic--Artin sector\\
Tate on a cubic fourfold; on K3 divisors or its square
 & published theorem & Does not itself assert the mixed maps (15.3), (15.10)\\
General fixed \(M(X)\) and arbitrary abelian comparison
 & named outside preprint + published theorem & van Hoften--Sempliner 2610.00259v2 + Milne 1999, Thm 7.1\\
CM rationality preprint 2610.00226v1; Milne note 2611.00244v1
 & not used & No rationality step imported from these here\\
A; M & not used & No analytic GRH input creates a cycle\\
\end{tabular}
\end{center}

Deligne's published absolute-Hodge theorem says that Hodge classes on
abelian varieties are absolute Hodge classes. It does not say that all
such classes are algebraic. Milne's published 1999 Theorem 7.1 derives
all-abelian Tate from algebraicity of Hodge classes on CM abelian
varieties. Replacing that algebraicity hypothesis by Deligne's
absolute-Hodge conclusion is not justified. The published substitutes
above work because their particular algebraic cycle classes are known,
not because ``absolute Hodge'' is redefined to mean ``algebraic''.

\paragraph{Reading in the language of the sheet.}
The Weil eigenvalues, their weights and norms control the spectral
object contemplated in \(Q'\) and in the boxed purity statement. A
forward Chow correspondence is the additional geometric comparison
from the variety to that object. Kunneth and Poincare duality turn that
comparison into a Tate class on a product, precisely (ARCH-15.3) or
(ARCH-15.10). In the geometric examples, graph, exceptional-divisor,
and Jacobian correspondences are the comparison itself; the Frobenius
polynomial alone is not being used in its place. This keeps separate
the sheet's spectral question and the cycle-theoretic content of its
conjecture forte.

\paragraph{Source locators.}
Exact primary URLs, inspected versions, and statement-level limits are
recorded in \texttt{cote15\_sources.md}. The core published locators are
Shioda--Katsura (1979), Lemmas 1.1--1.2, pp.~98--99, Remark 1.10,
p.~102; Milne (2001), Appendix A.7, author p.~21; Spie\ss\ (1999),
\emph{Math. Ann.} 314, 285--290; Zarhin (2004), Corollary 5.2 (author
preprint p.~8); Ito--Ito--Koshikawa (2021), Theorems 1.1 and 1.5,
pp.~2--3; Auel et al., Theorem 4.10, author-hosted journal PDF p.~20;
and Deligne (1982), Theorem 2.11. None of these bibliographic pointers
records an independent whole-proof verification of the cited published
theorem; the check ledger specifies what was actually read.
